The following records, documented in Chapter \ref{ch-network} are used by PFC:
\begin{itemize}
\item[] BUS
\item[] LINE
\item[] SWITCH
\item[] TRANSFO
\item[] TRFO
\item[] NRTP.
\end{itemize}

The additional (mandatory or optional) records only used in power flow computations are documented in this chapter.

\section{Load and shunt data}\label{sec-pf-bus}

Load and shunt data are attached to buses. They are specified in an extended version of the BUS record, as follows:

\hspace*{5ex}{\tt BUS NAME VNOM PLOAD QLOAD BSHUNT QSHUNT ;}

where:
\begin{itemize}
\item {\tt NAME} is the name of the bus. This is a string of at most 8 characters
\item {\tt VNOM} is the nominal voltage, in kV
\item {\tt PLOAD} is the total active power load at the bus, in MW. A positive value corresponds to power drawn from the network
\item {\tt QLOAD} is the total reactive power load at the bus, in Mvar. A positive value corresponds to power drawn from the network
\item {\tt BSHUNT} is the nominal reactive power of the shunt modelled as constant susceptance, in Mvar. This is the reactive power produced by the shunt under the nominal voltage of the bus, specified in its BUS record. BSHUNT is positive (resp. negative) for a shunt capacitor (resp. inductor)
\item {\tt QSHUNT} is the reactive power produced by the shunt modelled as constant power, in Mvar. QSHUNT is positive (resp. negative) for a shunt capacitor (resp. inductor).
\end{itemize}

The total reactive power $Q$ produced by the two shunt components is given, in Mvar, by:
\[ Q = BSHUNT \; (\frac{V}{V_{nom}})^2 + QSHUNT \]
where $V$ is the bus voltage and $V_{nom}$ the corresponding nominal voltage.

If no shunt is connected to the bus, BSHUNT and QSHUNT are set to zero.

If no load is connected to the bus, PLOAD and QLOAD are set to zero.

Let us recall that the PLOAD, QLOAD, BSHUNT and QSHUNT fields are ignored by RAMSES.


%----------------------------------------------------------------------------------
\section{Generator data}

Generators are specified by GENER records, as follows:

\hspace*{5ex}{\tt GENER NAME BUS P Q VIMP SNOM QMIN QMAX BR ;}

where:
\begin{itemize}
\item {\tt NAME} is the name of the generator. This is a string of a most 20 characters
\item {\tt BUS} is the name of the bus which the generator is connected to. This is a string of at most 8 characters
\item {\tt P} is the  active power produced by the generator, in MW
\item {\tt Q} is the reactive power produced by the generator, in Mvar. This value is ignored if the {\tt VIMP} field is nonzero
\item {\tt VIMP} is the voltage imposed by the generator, in pu. If {\tt VIMP} is zero, the bus is
treated as a PQ bus with the reactive power production set to {\tt Q}; if {\tt VIMP} is nonzero, the bus is treated as a PV bus and the {\tt Q} field is ignored
\item {\tt SNOM} is the nominal apparent power of the generator, in MVA
\item {\tt QMIN} is the lower reactive power limit, in Mvar. It is used only if {\tt VIMP} is nonzero
\item {\tt QMAX} is the upper reactive power limit, in Mvar. It is used only if {\tt VIMP} is nonzero
\item {\tt BR} is the on/off status of the generator breaker. A zero value indicates that the
breaker is open; any other value means that the breaker is closed.
\end{itemize}

For all generators with a nonzero value of {\tt VIMP}, the connection bus is initially assumed of the PV type.

If this entails exceeding the generator upper reactive power limit {\tt QMAX}, the bus switches to PQ type, the {\tt QMAX} limit is enforced, and Newton iterations continue. If subsequently the bus voltage gets larger than {\tt VIMP}, the bus switches back to PV type with the voltage imposed at the {\tt VIMP} value.

Similarly, if the generator lower reactive power limit {\tt QMIN} is exceeded, the bus switches to PQ type, the {\tt QMIN} limit is enforced, and Newton iterations continue. If subsequently the bus voltage gets smaller than {\tt VIMP}, the bus switches back to PV type with the voltage imposed at the {\tt VIMP} value.

Only one generator is allowed per bus.

All generators are memorized, even those which are disconnected. A disconnected generator is not involved  (it appears in the output results with a zero power output) but it can be put into service in the
dynamic simulation.

{\bf \large Remark 1}. For simplicity, a generator producing constant active and reactive powers can be modelled as a negative load using the BUS and no GENER record.

{\bf \large Remark 2}. There exists a variant of the GENER record with the following syntax:

\hspace*{5ex}{\tt GENER NAME BUS P Q V SNOM QMIN QMAX PMIN PMAX BR ;}

where {\tt PMIN} (resp. {\tt PMAX}) is the minimum (resp. maximum) active power that the generator can produce, in MW. If this record is present in the data, the {\tt PMIN} and {\tt PMAX} fields are ignored by STEPSS.


%----------------------------------------------------------------------------------
\section{Slack-bus specification}

The presence of a slack-bus is mandatory in power flow computations: indeed, not all buses can be of the PV or PQ type, since this would entail specifying the active power losses in the network, which are not known before performing the power flow calculation.

A generator of the PV type must be connected to the slack-bus. Its voltage magnitude, specified in its GENER record, is imposed at the bus, while the voltage phase angle is set to zero.

PFC can handle only one connected network (or island). If the network graph is not connected, {\sl \large only the connected sub-network including the slack-bus is treated}; the rest of the network is ignored with a warning message.

The SLACK record allows specifying which bus is the slack-bus. The syntax is as follows:

\hspace*{5ex}{\tt SLACK NAME ;}

where {\tt NAME} is the name of the bus. This is a string of at most 8 characters.

There must be exactly one SLACK record in the whole set of data.

%----------------------------------------------------------------------------------
\section{Static Var Compensators}

\subsection{Modelling}

The Static Var Compensator (SVC) is modelled as shown in Fig.~\ref{fig-pf-SVC}. $j$ is the {\sl \large monitored } bus, whose voltage is regulated, while $i$ is the {\sl \large controlled} bus, where the shunt susceptance $B$ is varied. $i$ and $j$ can be any two buses but usually $j$ is the transmission bus on the high-voltage side of the SVC step-up transformer, while the SVC is connected to bus $i$.

\begin{figure}[h!]
\centering
\includegraphics[scale=0.8]{SVC.pdf}
\caption{Block-diagram of the static var compensator}\label{fig-pf-SVC}
\end{figure}

The SVC being assumed lossless, the active current injected at bus $i$ is zero ($I_{Pi} = 0$), while the reactive current takes on one of the following forms:
\begin{eqnarray}
I_{Qi} & = & B V_i = G (V_j^o - V_j) V_i \;\;\;\; \mbox{under voltage control} \label{eq-svc-vcontrol} \\
I_{Qi} & = & B_{max} V_i \;\;\;\; \mbox{under upper susceptance limit} \label{eq-svc-Bmax}\\
I_{Qi} & = & B_{min} V_i \;\;\;\; \mbox{under lower susceptance limit.} \label{eq-svc-Bmin}
\end{eqnarray}

Although reference is made to an SVC, the model can be used in general for a component controlling voltage with a droop.

\subsection{Data format}

The record, with keyword SVC, includes the following fields:

\hspace*{5ex}{\tt SVC NAME CON\_BUS MON\_BUS V0 Q0 SNOM BMAX BMIN G BR ;}

where:
\begin{itemize}
\item {\tt NAME} is the name of the compensator. This is a string of a most 20 characters
\item {\tt CON\_BUS} is the name of the controlled bus. This is a string of a most 8 characters
\item {\tt MON\_BUS} is the name of the monitored bus. This is a string of a most 8 characters
\item {\tt V0} is the voltage setpoint $V_j^o$, in pu (see Fig.~\ref{fig-pf-SVC}) 
\item {\tt Q0} is the reactive power setpoint, in Mvar. If {\tt V0} is set to zero, the SVC is treated under constant power, with $P = 0$ and $Q = Q0$, and no limit is
tested. If {\tt V0} is nonzero, the {\tt Q0} field is ignored
\item {\tt SNOM} is the nominal reactive power of the SVC, in Mvar
\item {\tt BMAX} is the maximal nominal reactive power, in Mvar. This is the reactive 
power produced by the compensator under $V_i = 1$~pu, when $B$ is at
the maximal value $B_{max}$.
\item {\tt BMIN} is the minimal nominal reactive power, in Mvar. This is the reactive
produced by the compensator under $V_i = 1$~pu, when $B$ is at
the minimal value $B_{min}$.
\item {\tt G} is the gain $G$, in pu on the ($V_B$, {\tt SNOM}) base, where $V_B$ is the nominal voltage at the controlled bus, as given by its BUS record
\item {\tt BR} is the on/off status of the SVC breaker. A zero value indicates that
the breaker is open, any other value means that it is closed.
\end{itemize}

It is common for {\tt BMAX} to be positive and {\tt BMIN} negative but other combinations are allowed.

For all SVCs with a nonzero value of {\tt V0}, Eq.~(\ref{eq-svc-vcontrol}) is solved initially.

If this entails exceeding the susceptance upper reactive power limit {\tt BMAX}, Eq.~(\ref{eq-svc-Bmax}) is substituted, the {\tt BMAX} limit is enforced, and Newton iterations continue. If subsequently $G(V_j^o-V_j)< B_{max}$, the program reverts to Eq.~\ref{eq-svc-vcontrol} and keeps on iterating.

Similarly, if the SVC lower suscepance limit {\tt BMIN} is exceeded, Eq.~(\ref{eq-svc-Bmin}) is substituted, the {\tt BMIN} limit is enforced, and Newton iterations continue. If subsequently $G(V_j^o-V_j) > B_{min}$, the program reverts to Eq.~(\ref{eq-svc-vcontrol}) and keeps on iterating.

Only one SVC is allowed per bus. 

It is not allowed to connect both a generator and an SVC to the same bus.

All SVCs are memorized, even those which are disconnected. A disconnected SVC is not involved (it appears in the output results with a zero power output) but it can be put into service in the
dynamic simulation.


%----------------------------------------------------------------------------------
\section{Transformer ratio adjustment for voltage control}\label{sec-pf-ltc}

\subsection{Modelling}

PFC can adjust the ratio of a designated transformer with the objective of bringing a controlled voltage inside a deadband $[V_{des}-\epsilon \;\; V_{des}+\epsilon]$, where $V_{des}$ is the desired voltage and $\epsilon$ is a tolerance.

The ratio is changed in discrete steps (as in the real life), between a minimum and a maximum value. During the computation, the ratio is changed by one step at a time, after which Newton iterations are performed until convergence is achieved. The process is repeated until the controlled voltage falls in the deadband. When multiple transformers are adjusted, some may reach their deadbands before others.

\subsection{First data format}

The first way to specify ratio adjustement is through the TRFO record. The following refers to the material  presented in Section~\ref{sec-netw-trfo}.

The controlled bus is {\tt CONBUS}. This must be one of the two ending buses of the transformer. An empty or blank string {\sl \large enclosed within quotes} indicates that the transformer ratio is not to be adjusted. In this case, however, dummy values must be provided for each of the fields listed below.

The fields of concern are:

\begin{itemize}
\item {\tt NFIRST}: the ratio in \% corresponding to the first tap position. This is the lower bound on the transformer ratio
\item {\tt NLAST}: the ratio in \% corresponding to the last tap position. This is the upper bound on the transformer ratio
\item {\tt NBPOS}: the total number of tap positions including the first and the last
\item {\tt TOLV}: the voltage tolerance $\epsilon$, in pu
\item {\tt VDES}: the desired voltage $V_{des}$, in pu.
\end{itemize}

The transformer ratio $n$ corresponding to position $p$ ($1 \le p \le {\tt NBPOS}$) of the tap changer is given by:
\begin{equation}\label{equ-pf-nvalue}
n = \frac{{\tt NFIRST}}{100} + \frac{p-1}{{\tt NBPOS}-1} \; \frac{{\tt NLAST} - {\tt NFIRST}}{100}
\end{equation}
The value of $p$ is increased or decreased by one at each adjustment iteration.

An initial value of the transformer ratio is given by the {\tt N} field of the TRFO record. If the transformer is to be adjusted, before starting the power flow computation that value is adjusted to coincide with the nearest tap position, in accordance with Eq.~(\ref{equ-pf-nvalue}).

\subsection{Second data format}

The second way to specify ratio adjustement is through a separate record with keyword LTC-V. The syntax is as follows:

\hspace*{5ex}{\tt LTC-V NAME CON\_BUS NFIRST NLAST NBPOS TOLV VDES ;}

where {\tt NAME} is the name of the controlled transformer (a string of at most 20 characters) and all the other fields have the same meaning as for the TRFO record.

A transformer can be controlled by a single tap changer only. It is more natural to use the LTC-V record in association with a TRANSFO record. Hovever, the LTC-V record can be associated with a TRFO record, provided that no adjustment is specified in the latter. 



%----------------------------------------------------------------------------------
\section{Phase-shifting transformer ratio adjustment for power control}

\subsection{Modelling}

PFC can also adjust the phase angle of a transformer with the objective of bringing the active power flow in a branch inside a deadband $[P_{des}-\epsilon \;\; P_{des}+\epsilon]$, where $P_{des}$ is the desired power flow and $\epsilon$ is a tolerance.

The adjustment is very similar to that of in-phase transformers for voltage control detailed in  Section~\ref{sec-pf-ltc}.

\subsection{Data format}

The record, with keyword PSHIFT-P, includes the following fields:

\hspace*{5ex}{\tt PSHIFT-P CONTRFO MONBRANCH PHAFIRST PHALAST NBPOS SIGN PDES TOLP ;}

where:
\begin{itemize}
\item {\tt CONTRFO} is the name for the transformer whose phase angle is to be adjusted. This is a string of at most 20 characters, defined in either a TRFO or a TRANSFO record. If the transformer does
not exist, the whole record is ignored with a warning message
\item {\tt MONBRANCH} is the name of the branch in which the active power flow $P$ is
monitored. This is a string of at most 20 characters, defined in either a LINE, a TRFO or a TRANSFO record.  The sign convention is the following: $P$ is the active power flow leaving the first bus specified in the LINE, TRFO or TRANSFO record and entering the branch of concern
\item {\tt PHAFIRST} is the phase angle $\phi$, in degrees, corresponding to the first tap position. This is the lower bound on $\phi$
\item {\tt PHALAST} is the phase angle $\phi$, in degrees, corresponding to the last tap position. This is the upper bound on $\phi$
\item {\tt NBPOS} is the number of tap positions
\item {\tt SIGN} is an indication of the direction in which the phase angle $\phi$ must be adjusted to reach the objective. A value of $1$ indicates that $\phi$ must be increased to increase the active power flow in the monitored branch. A value of $-1$ indicates that it must be decreased. Any other value is invalid and causes the program to stop
\item {\tt PDES} is the desired active power flow, in MW
\item {\tt TOLP} is the tolerance $\epsilon$, in MW.
\end{itemize}

The phase angle $\phi$ corresponding to position $p$ ($1 \le p \le {\tt NBPOS}$) of the tap changer is given by:
\begin{equation}\label{equ-pf-phivalue}
\phi = {\tt PHAFIRST} + \frac{p-1}{{\tt NBPOS}-1} \; ({\tt PHALAST} - {\tt PHAFIRST})
\end{equation}
The value of $p$ is increased or decreased by one at each adjustment iteration.

PFC performs a sensitivity analysis to determine whether the phase angle should be increased or decreased. If this analysis indicates a direction opposite to what is specified in {\tt SIGN}, a warning is issued and the value of {\tt SIGN} is ignored. On output, when it produces a file with the records updated, PFC sets {\tt SIGN} to the value corresponding to its sensitivity analysis.

An initial value of the phase angle is given by the {\tt PHI} field of the TRANSFO record. If the transformer is to be adjusted, before starting the power flow computation, that value is adjusted to coincide with the nearest tap position, in accordance with Eq.~(\ref{equ-pf-phivalue}).

A transformer cannot be controlled by both an LTC-V and a PSHIFT-P record.

Only one PSHIFT-P record per transformer is allowed. The PSHIFT-P record is intended to be used in association with a TRANSFO record. However, it is allowed to associate it with a TRFO record in spite of the fact that the latter assumes a zero phase angle. In this case, the angle will be initialized to zero and will be controlled as specified in the PSHIFT-P record.


%----------------------------------------------------------------------------------
\section{Bus voltages: initial values and results}\label{sec-lfresv}

On output, PFC produces a file with the computed bus voltage magnitudes and phase angles. The latter are specified in records with keyword LFRESV. The syntax is as follows:

\hspace*{5ex}{\tt LFRESV BUS MODV PHASV ;}

where:
\begin{itemize}
\item {\tt BUS} is the name of the bus. This is a string of a most 8 characters defined in a BUS record
\item {\tt MODV} is the bus voltage magnitude, in pu
\item {\tt PHAV} is the bus voltage phase angle, in radian, the slack bus being the reference.
\end{itemize}

If LFRESV records are provided on input, they are used as initial voltages of the Newton iterations, at the specified buses. 

If no LFRESV record is specified at a bus:
\begin{itemize}
\item the voltage magnitude is initialized to 1 pu if the bus is of the PQ type or to the value imposed by the generator in case of a PV bus
\item the phase angle is initialized to 0 degree.
\end{itemize}

Thus, if the LFRESV records obtained from a first run of PFC are added to the data files, the other data being unchanged, no Newton iteration is going to be performed since we are already at the solution. This is an easy way to check that system data come with their corresponding voltages.


%----------------------------------------------------------------------------------
\section{Share of records beween PFC and RAMSES}

A summary of the records used by respectively PFC and RAMSES is given in Table~\ref{tab-records}.

\begin{table}[h!]
\centering
\caption{records used by PFC and RAMSES, respectively}\label{tab-records}
\begin{tabular}{|c|c|c|}
\hline
record & in PFC & in RAMSES \\ \hline \hline
BUS & all 6 fields used & only first two fields used \\ \hline
LINE & used & used \\ \hline
SWITCH & used & used \\ \hline
NRTP & used & used \\ \hline
TRANSFO & used & used \\ \hline
TRFO & used & only fields 1 to 9 and 15 used \\ \hline
SHUNT & ignored & used \\ \hline
GENER &  used & ignored \\ \hline
SVC & used & ignored \\ \hline
SLACK & used & used \\ \hline
LFRESV & used & used \\ \hline
LTC-V & used & ignored \\ \hline
PSHIFT-P & used & ignored \\ \hline
\end{tabular}
\end{table}


%----------------------------------------------------------------------------------
\section{Computation control parameters}\label{sec-pfc-ctl}

\subsection{Parameters}

PFC performs Newton(-Raphson) iterations to solve the power flow equations. At the $k$-th iteration, the following indices are computed:
\begin{itemize}
\item[] $\displaystyle \epsilon_P = \max_i |f_i({\bf v}^{(k)}, {\bm \theta}^{(k)}) - P^o_i|$ the largest absolute mismatch of the active power equations
\item[] $\displaystyle \epsilon_Q = \max_i |g_i({\bf v}^{(k)}, {\bm \theta}^{(k)}) - Q^o_i|$ the largest absolute mismatch of the reactive power equations
\item[] $\displaystyle \epsilon_S = \max_i \sqrt{(f_i({\bf v}^{(k)}, {\bm \theta}^{(k)}) - P^o_i)^2 + 
    (g_i({\bf v}^{(k)}, {\bm \theta}^{(k)}) - Q^o_i)^2}$ the largest apparent power mismatch.
\end{itemize}
Convergence is achieved once :
\begin{itemize}
\item both $\epsilon_P$ and $\epsilon_Q$ are below specified thresholds
\item all controls on transformer ratios and phase shifts are satisfied, and
\item all generators and SVCs are within their reactive limits.
\end{itemize}

$\epsilon_S$ is used to check that the solution is accurate enough to :
\begin{itemize}
\item ``freeze" the Jacobian matrix (in order to save some computing time)
\item check the generator and SVC reactive power limits and enforce the latter if needed
\item adjust the transformer ratios and phase shifs to control voltage magnitudes and active power flows, if specified in the data.
\end{itemize}

Finally, the iterations stop as soon as divergence is detected. To this purpose the quadratic index:
\[  \varphi(k) = \sum_i \sqrt{(f_i({\bf v}^{(k)}, {\bm \theta}^{(k)}) - P^o_i)^2 + (g_i({\bf v}^{(k)}, {\bm \theta}^{(k)}) - Q^o_i)^2} \]
is monitored. Under normal convergence conditions, $\varphi$ decreases from one iteration to the next. Therefore, the increase in $\varphi$ is used to detect divergence. The algorithm stops at the iteration $k$ such that :
\[ \varphi (k) > 1.1 \: \varphi(k-1) \]
However, this test is skipped at any iteration $k$ that follows the switching of generators or SVCs under limit, or the adjustment of transformer ratios and phase shifts\footnote{indeed, these adjustments cause increases in $\varphi$ that have nothing to do with divergence}.

\subsection{Records}

The records detailed herafter are used to control the computation. They all start with a \$ to distinguish them from the other records. They all have a single field, as detailed in Table~\ref{tab-control-param}. The default value assigned in the absence of the record is given in the fourth column.

\begin{table}[h!]
\centering
\caption{Computation control parameters}\label{tab-control-param}
\begin{tabular}{|c|c|c|c|}
\hline
record & meaning of field & unit & default value \\
\hline \hline
\tt \$SBASE   & base power  & MVA & 100 \\
              & (on which pu values are expressed) & & \\ \hline
\tt \$TOLAC   & $\epsilon_P$ & MW & 0.1 \\ \hline
\tt \$TOLREAC & $\epsilon_Q$ & Mvar & 0.1 \\ \hline
\tt \$NBITMA  & max number of iterations & - & 20 \\ \hline
\tt \$MISQLIM & value of $\epsilon_S$ below which the reactive & MVA & 20 \\
              & power limits are checked and enforced. &     & \\ 
              & Set to 0 to skip the limit check & & \\ \hline
\tt \$MISBLOC & value of $\epsilon_S$ below which & MVA & 10 \\ 
              & the Jabobian is kept constant. & & \\ \hline

\tt \$MISADJ  & value of $\epsilon_S$ below which & MVA & 10 \\
              & transformers are adjusted. & & \\ 
              & Set to 0 to skip adjustment & & \\ \hline
\tt \$DIVDET  & set to 1 to activate divergence & - & 0 \\
              & test during iterations.  & & \\ 
              & Set to 0 to skip test & & \\ \hline
\end{tabular}
\end{table}

\clearpage

